\chapter{System Overview, Technology, and Structure}

\section{What Is This System?}

The Workforce Management Platform is a server-rendered web application
that helps a single organization answer the practical, day-to-day
questions a business asks about its hourly workforce:

\begin{itemize}[nosep]
  \item Who is scheduled to work, and who actually worked?
  \item How many hours did each employee work, and how much of that was
        overtime?
  \item Who is on approved leave, and who has a leave request waiting?
  \item How much is labor costing the organization, per department?
  \item What privileged changes has anyone made recently?
\end{itemize}

The system is built around six core domain concepts, each with its own
models and service module:

\begin{center}
\begin{tikzpicture}[node distance=0.5cm]
  \node[flownode, text width=3.1cm] (dept) {Departments};
  \node[flownode, text width=3.1cm, right=of dept] (emp) {Employees};
  \node[flownode, text width=3.1cm, right=of emp] (sched) {Scheduling\\(Shifts)};
  \node[flownode, text width=3.1cm, below=of dept] (att) {Attendance};
  \node[flownode, text width=3.1cm, below=of emp] (leave) {Leave};
  \node[flownode, text width=3.1cm, below=of sched] (cost) {Overtime \&\\Labor Cost};
\end{tikzpicture}
\end{center}

Two of these concepts are deliberately kept separate and are never
conflated anywhere in the codebase:

\begin{itemize}
  \item \textbf{Scheduling} (the \code{Shift} model) represents
        \emph{planned} work --- what an employee is expected to work.
  \item \textbf{Attendance} (the \code{AttendanceEntry} model)
        represents \emph{actual} work --- what really happened, recorded
        via clock-in/clock-out.
\end{itemize}

Working hours, overtime, and labor cost are all \emph{derived} from
these two record types rather than stored as independent numbers -- see
Chapter~4 for exactly how.

\subsection*{Who Uses It}

Three roles share the same application, each seeing a different slice
of it, enforced by \code{AccessScope} (Chapter~2):

\begin{itemize}[nosep]
  \item \textbf{Admin} --- full access to one organization: every
        employee, every department, pay rates, labor cost detail, the
        audit log, and every management action.
  \item \textbf{Manager} --- restricted to the departments they are
        assigned to manage (\code{department\_managers}). Can manage
        schedules, attendance, and leave for their own people, and see
        department-level cost \emph{totals} --- never an individual
        employee's pay rate or exact cost.
  \item \textbf{Employee} --- self-service only: their own schedule,
        their own attendance (clock in/out), their own leave requests.
\end{itemize}

\section{Product Scope}

The table below reflects what actually exists in the codebase today,
verified against the route and service modules in \file{app/routes/}
and \file{app/services/}.

\begin{longtable}{@{}p{4.2cm}p{2.7cm}p{7.3cm}@{}}
\toprule
\textbf{Capability} & \textbf{Status} & \textbf{Description} \\
\midrule
\endhead
Authentication & Implemented & Email/password login, Argon2 hashing,
  account lockout after repeated failures, absolute session expiry. \\
User accounts \& roles & Implemented & \code{admin} / \code{manager} /
  \code{employee}, enforced by DB \code{CHECK} constraint and
  \code{AccessScope}. \\
Departments & Implemented & CRUD (admin), soft-deactivation, manager
  assignment via \code{department\_managers}. \\
Employees & Implemented & CRUD, employment status lifecycle
  (active/inactive/terminated), department assignment. \\
Scheduling (shifts) & Implemented & Draft/published/cancelled shifts,
  employee assignment, overlap and leave-conflict prevention. \\
Check-in / check-out & Implemented & Self-service and admin/manager
  on-behalf-of clock-in/out, single-open-entry enforcement. \\
Breaks & \PARTIAL & A single \code{break\_minutes} integer per shift/
  entry only --- there is no separate break start/end event, no
  multiple-breaks-per-shift model, and no \code{BreakSession} table
  anywhere in the schema. \\
Attendance correction & Implemented & Admin/manager correction with a
  mandatory reason; stale open entries auto-flagged \code{needs\_review}. \\
Working hours & Implemented & Derived from closed attendance entries
  only; scheduled-vs-worked comparison per employee/day. \\
Overtime & Implemented & Configurable daily/weekly tiered policy,
  pure-function tiering, no double-counting across tiers. \\
Leave & Implemented & Request/approve/reject/cancel lifecycle, overlap
  prevention, shift-conflict detection on approval. \\
Leave balances / accrual & \PLANNED & \code{LeaveType} is a catalog
  only; no "days remaining" ledger exists. \\
Labor cost & Implemented & Decimal, per-line-item rounded cost;
  department totals (admin/manager) and per-employee breakdown (admin
  only). \\
Paid leave costing & \PLANNED & \code{LeaveType.is\_paid} is stored but
  never read by the labor-cost module; approved leave currently
  contributes \$0 to every cost figure (documented explicitly in
  \file{app/services/labor\_cost.py}). \\
Reports & Implemented & Overtime-hours report and hours-trend report,
  both department-scoped. \\
Audit log & Implemented & Append-only, admin-only, paginated,
  date-range filtered; a lighter "Recent Activity" view over the same
  data. \\
Reduced information density: workforce demand forecasting, AI-assisted
scheduling, schedule optimization & \PLANNED & Explicitly out of scope
  for MVP~1 per \file{CLAUDE.md}. \\
\bottomrule
\end{longtable}

\section{Technology Stack}

Every entry below is a real, currently-used dependency, verified
against \file{requirements.txt} and the import statements in
\file{app/}.

\subsection*{Python 3.14}
\textbf{What:} The application's implementation language.
\textbf{Why:} Type hints, dataclasses, and the standard \code{zoneinfo}
module (used throughout for timezone-aware date arithmetic) are used
extensively across the service layer.
\textbf{Where:} Everywhere under \file{app/}.

\subsection*{Flask 3.1}
\textbf{What:} A minimal WSGI web framework.
\textbf{Why:} Matches the project's "modular monolith, not a
microservice framework" architectural preference (\file{CLAUDE.md}) ---
Flask imposes very little structure of its own, leaving the
routes/services/models layering to be an explicit project convention
rather than a framework-enforced one.
\textbf{Where:} \code{create\_app()} in \file{app/\_\_init\_\_.py};
every blueprint in \file{app/routes/}.
\textbf{Responsibility:} HTTP request/response handling, routing,
session cookies, template rendering.

\subsection*{SQLAlchemy 2.0 (ORM) + PostgreSQL}
\textbf{What:} SQLAlchemy is the Object-Relational Mapper; PostgreSQL is
the relational database.
\textbf{Why:} The domain has many real relational invariants (no
overlapping shifts, no double clock-in, effective-dated pay rates and
overtime policies) that are enforced with PostgreSQL-specific features
--- \code{EXCLUDE} constraints using \code{gist} indexes over
\code{tstzrange}/\code{daterange}, partial unique indexes, and
composite foreign keys --- not just application code.
\textbf{Where:} \file{app/models/*.py} (declarative models, using the
2.0-style \code{Mapped[...]} annotations); every \code{db.session}
query in \file{app/services/*.py}.
\textbf{Responsibility:} Schema definition, query construction,
transaction management.

\subsection*{Flask-Migrate (Alembic)}
\textbf{What:} Database migration tooling built on Alembic.
\textbf{Why:} Schema changes must go through versioned, reviewable
migration scripts rather than ad hoc DDL (\file{.claude/rules/database.md}).
\textbf{Where:} \file{migrations/versions/0001} through \file{0016}.
\textbf{Responsibility:} Applying and reversing schema changes in a
controlled, ordered sequence.

\subsection*{Flask-Login}
\textbf{What:} Session-based authentication helper.
\textbf{Why:} Provides \code{current\_user}, the \code{login\_required}
decorator, and the user-loader hook this project uses to invalidate a
session the instant an account is deactivated or locked (see
Chapter~2).
\textbf{Where:} \code{app.extensions.login\_manager};
\code{app/models/user.py}'s \code{load\_user}; \file{app/auth/decorators.py}.

\subsection*{Flask-WTF}
\textbf{What:} WTForms integration with built-in CSRF protection.
\textbf{Why:} Every state-changing form in the application needs CSRF
protection and server-side field validation; Flask-WTF provides both
with one consistent pattern.
\textbf{Where:} \code{app.extensions.csrf}; every form class in
\file{app/forms.py}.

\subsection*{argon2-cffi}
\textbf{What:} A Python binding for the Argon2 password-hashing
algorithm.
\textbf{Why:} Argon2 is the current recommended password hash for new
applications (memory-hard, resistant to GPU cracking), used here
instead of an older algorithm like bcrypt or PBKDF2.
\textbf{Where:} \file{app/auth/passwords.py} (\code{hash\_password} /
\code{verify\_password}), wrapping \code{argon2.PasswordHasher}.

\subsection*{Gunicorn}
\textbf{What:} A production-grade WSGI HTTP server for Python.
\textbf{Why:} Flask's own development server is explicitly
unsuitable for production; Gunicorn is the process that actually serves
requests in the Docker image.
\textbf{Where:} \file{Dockerfile}'s \code{CMD}: \code{gunicorn -b
0.0.0.0:8000 wsgi:app}.

\subsection*{HTML, CSS, and Vanilla JavaScript}
\textbf{What:} Server-rendered Jinja2 templates, a hand-written CSS
design-token system, and a small amount of unframeworked JavaScript.
\textbf{Why:} \file{CLAUDE.md} explicitly forbids introducing a
frontend framework (React/Vue/Angular) without justification; the UI
is rendered entirely server-side.
\textbf{Where:} \file{templates/}, \file{static/css/}, \file{static/js/main.js}.
See Chapter~6 for the design-system details.

\subsection*{pytest}
\textbf{What:} The test framework.
\textbf{Why:} Used for unit, integration, and route-level tests across
the whole codebase (458 tests at the time of writing).
\textbf{Where:} \file{tests/}. See Chapter~7.

\section{High-Level Architecture}

The application is a \textbf{modular monolith}: a single deployable
Flask process, internally organized into clearly separated layers,
rather than a collection of independently deployed services. This is a
deliberate choice (\file{.claude/rules/architecture.md}) --- the domain
is not large enough to justify the operational cost of microservices,
and a monolith keeps transactional integrity simple (a leave approval
that also writes an audit entry is one database transaction, not a
distributed one).

\begin{center}
\begin{tikzpicture}[node distance=0.5cm]
  \node[flownode] (browser) {Browser (HTML forms, no JS framework)};
  \node[flownode, below=of browser] (flask) {Flask Application (\code{create\_app})};
  \node[smallnode, below left=0.6cm and 0.4cm of flask] (routes) {Routes\\ (\file{app/routes/})};
  \node[smallnode, below right=0.6cm and 0.4cm of flask] (auth) {Authentication \&\\ Authorization\\ (\file{app/auth/})};
  \node[flownode, below=1.6cm of flask] (services) {Services (business logic) --- \file{app/services/}};
  \node[flownode, below=of services] (orm) {SQLAlchemy ORM --- \file{app/models/}};
  \node[flownode, below=of orm] (pg) {PostgreSQL};

  \draw[flowarrow] (browser) -- (flask);
  \draw[flowarrow] (flask) -- (routes);
  \draw[flowarrow] (flask) -- (auth);
  \draw[flowarrow] (routes.south) -- (services.north west);
  \draw[flowarrow] (auth.south) -- (services.north east);
  \draw[flowarrow] (services) -- (orm);
  \draw[flowarrow] (orm) -- (pg);
\end{tikzpicture}
\end{center}

\subsection*{Why Routes Stay Thin}

Every route module's docstring states the same rule: a route builds an
\code{AccessScope} from the signed-in user and delegates \emph{all}
authorization and data access to a service function --- no route
queries the database directly (the one historical exception,
\code{app.routes.audit}, was fixed to go through
\code{app.services.audit.list\_entries} instead, per that module's own
``Round C fix'' docstring). This keeps HTTP concerns (form parsing,
flashing messages, redirects) separate from business rules, and means a
business rule is only ever expressed once.

\subsection*{Why Business Logic Belongs in Services}

Every function in \file{app/services/*.py} takes the caller's
\code{AccessScope} as its first argument and re-checks authorization
itself, \emph{independent of whatever the route layer already checked}.
This "belt-and-suspenders" pattern means that even if a route were
missing a decorator, the service function underneath would still refuse
an unauthorized call. It also means services are usable from multiple
entry points (a route, a CLI command, a future API) without
re-implementing authorization each time.

\subsection*{Why Models Represent Persistence Only}

Models in \file{app/models/} are declarative SQLAlchemy classes with
constraints and relationships, and carry essentially no business logic
of their own (a few carry a \code{\_\_repr\_\_} for debugging and
nothing else). Validation and business rules live in the service layer;
the database's own constraints (Chapter~3) are the last line of
defense, not the primary one.

\subsection*{How Authorization Is Enforced}

See Chapter~2 for the full authorization model. In summary: a route
decorator (\code{login\_required} / \code{role\_required}) provides a
first, coarse check; every service function re-checks the caller's
role and organizational/department scope before touching data; and
several database constraints (composite foreign keys tying a child row's
organization to its parent's) make a cross-tenant write impossible even
if every application-level check were somehow bypassed.

\section{Complete Project Tree}

The tree below is the actual top-level layout of the repository (test
cache, \file{.git}, and Python bytecode caches omitted), matching
\file{.claude/rules/project-structure.md}'s canonical layout.

\begin{lstlisting}[basicstyle=\ttfamily\footnotesize,frame=single]
job/
|-- app/
|   |-- auth/            Authentication & authorization helpers
|   |   |-- decorators.py    login_required / role_required
|   |   |-- passwords.py     Argon2 hashing wrapper
|   |   |-- redirects.py     Open-redirect-safe "next" handling
|   |   |-- scope.py         AccessScope + get_scoped_or_404
|   |   `-- service.py       Login business logic, lockout policy
|   |-- models/          SQLAlchemy declarative models (13 tables)
|   |-- routes/          Flask blueprints (thin HTTP layer)
|   |-- services/        Business logic (the bulk of the application)
|   |-- cli.py           Flask CLI commands (attendance flag-stale, seed)
|   |-- config.py        Environment-driven configuration classes
|   |-- errors.py        400/403/404/413/500 error handlers
|   |-- extensions.py    Flask extension singletons (db, login, csrf...)
|   |-- forms.py         WTForms form classes
|   `-- __init__.py      create_app() application factory
|-- migrations/
|   `-- versions/        16 ordered Alembic migration scripts
|-- templates/           Jinja2 templates (one subfolder per blueprint)
|   |-- components/      Shared macros: icons, badges, sidebar, topbar
|   `-- errors/          400/403/404/413/500 error pages
|-- static/
|   |-- css/             tokens.css, base.css, components.css, layout.css
|   |-- fonts/           Self-hosted variable fonts (woff2)
|   `-- js/main.js       Small vanilla-JS behavior layer
|-- tests/               pytest suite (458 tests)
|   |-- unit/            Pure-function tests, no DB/Flask app
|   |-- integration/     Service-layer tests against a real test DB
|   `-- routes/          Full HTTP request/response tests
|-- docs/                This documentation (see docs/system-documentation/)
|-- CLAUDE.md            Project-wide AI/developer instructions
|-- .claude/rules/       Per-topic project rules (architecture, security, ...)
|-- docker-compose.yml   app + postgres services for local/containerized use
|-- Dockerfile           Production image: Gunicorn serving wsgi:app
|-- requirements.txt     Pinned Python dependencies
`-- wsgi.py              WSGI entrypoint (create_app())
\end{lstlisting}

\subsection*{Important Directories}

\textbf{\file{app/routes/}} --- One blueprint per domain area
(attendance, audit, auth, dashboard, departments, employees,
labor\_cost, leave, main, schedule). Each file's own docstring restates
the "thin route, delegate to services" rule.

\textbf{\file{app/services/}} --- The actual business logic: attendance,
audit, availability, departments, employees, labor\_cost, leave,
overtime, pay\_rates, reports, scheduling, working\_hours. This is the
directory a new developer should read first to understand what the
system actually does (see Chapter~9's "How to add a feature").

\textbf{\file{app/auth/}} --- Everything authorization-related lives
here, isolated from both routes and services so it can be reasoned
about (and reviewed) as a single unit. \code{AccessScope} in
\code{scope.py} is the single data structure every authorization
decision in the codebase is built from.

\textbf{\file{migrations/versions/}} --- The authoritative history of
the schema. Because Alembic migrations are additive scripts, the actual
deployed schema is the cumulative effect of all 16 files, not just the
current state of \file{app/models/} (see Chapter~3, \S3.4, on why the
two are cross-checked rather than assumed identical).

\textbf{\file{tests/}} --- Organized by test \emph{kind}, not by
feature: \file{unit/} needs no database or Flask app context at all
(e.g.\ the overtime tiering math); \file{integration/} exercises a
service function against a real PostgreSQL test database; \file{routes/}
drives the application exactly as an HTTP client would, including
session cookies and CSRF tokens.
